\documentclass{fphw}

\title{Problem Set 3: Rotational Dynamics}
\author{Jamie Rivera}
\class{Classical Mechanics I}
\professor{Dr. A. Okafor}
\institute{Department of Physics}
\date{\today}

\usepackage{amsmath}
\usepackage{amssymb}
\usepackage{graphicx}

\begin{document}

\maketitle

\begin{problem}
A uniform solid disk of mass $M$ and radius $R$ rolls without slipping down
an incline of angle $\theta$. Find the linear acceleration of the disk's
center of mass.
\end{problem}

\section*{Solution}

The disk has moment of inertia $I = \tfrac{1}{2}MR^2$ about its center. With
rolling without slipping, $a = R\alpha$. Newton's second law along the
incline and the rotational equation about the center give:
\begin{align*}
	Mg\sin\theta - f &= Ma, \\
	fR &= I\alpha = \tfrac{1}{2}MR^2 \cdot \frac{a}{R}.
\end{align*}
Solving the pair for $a$:
\[
	a = \frac{2}{3}g\sin\theta.
\]

\begin{problem}
Two blocks of mass $m_1$ and $m_2$ are connected by a massless string over a
frictionless, massless pulley. Derive the acceleration of the system and the
tension in the string.
\end{problem}

\section*{Solution}

Let $m_1 > m_2$, so $m_1$ descends. Newton's second law for each block:
\begin{align*}
	m_1 g - T &= m_1 a, \\
	T - m_2 g &= m_2 a.
\end{align*}
Adding the two equations eliminates $T$:
\[
	a = \frac{(m_1 - m_2)g}{m_1 + m_2}, \qquad
	T = \frac{2 m_1 m_2 g}{m_1 + m_2}.
\]

\section{Conceptual questions}

\begin{enumerate}
	\item Why does the rolling disk in Problem 1 accelerate more slowly than
	a block sliding down the same frictionless incline ($a = g\sin\theta$)?
	\item What happens to the tension $T$ in Problem 2 in the limit
	$m_1 = m_2$? Does that match your physical intuition?
\end{enumerate}

\end{document}
